% =====================================================================
%  Estilo visual del Aula Interactiva, llevado al papel.
%
%  Reproduce el lenguaje de la herramienta: los bloques temáticos con su
%  color, y los mismos componentes narrativos (Definición, Mini-historia,
%  Trampa, Fórmula anotada, Hilo, Cierre, Puente). La idea es que un
%  estudiante que viene de la pantalla reconozca las mismas señales en la
%  hoja impresa, sin tener que aprender un segundo sistema.
%
%  Interruptores, que el documento fija DESPUÉS de % =====================================================================
%  Estilo visual del Aula Interactiva, llevado al papel.
%
%  Reproduce el lenguaje de la herramienta: los bloques temáticos con su
%  color, y los mismos componentes narrativos (Definición, Mini-historia,
%  Trampa, Fórmula anotada, Hilo, Cierre, Puente). La idea es que un
%  estudiante que viene de la pantalla reconozca las mismas señales en la
%  hoja impresa, sin tener que aprender un segundo sistema.
%
%  Interruptores, que el documento fija DESPUÉS de % =====================================================================
%  Estilo visual del Aula Interactiva, llevado al papel.
%
%  Reproduce el lenguaje de la herramienta: los bloques temáticos con su
%  color, y los mismos componentes narrativos (Definición, Mini-historia,
%  Trampa, Fórmula anotada, Hilo, Cierre, Puente). La idea es que un
%  estudiante que viene de la pantalla reconozca las mismas señales en la
%  hoja impresa, sin tener que aprender un segundo sistema.
%
%  Interruptores, que el documento fija DESPUÉS de % =====================================================================
%  Estilo visual del Aula Interactiva, llevado al papel.
%
%  Reproduce el lenguaje de la herramienta: los bloques temáticos con su
%  color, y los mismos componentes narrativos (Definición, Mini-historia,
%  Trampa, Fórmula anotada, Hilo, Cierre, Puente). La idea es que un
%  estudiante que viene de la pantalla reconozca las mismas señales en la
%  hoja impresa, sin tener que aprender un segundo sistema.
%
%  Interruptores, que el documento fija DESPUÉS de \input{estilo-aula}:
%     \docentetrue / \docentefalse    la solución paso a paso
%     \espaciotrue / \espaciofalse    espacio en blanco para resolver
%
%  Se declaran acá y no en cada documento: si el documento hace su propio
%  \newif antes del \input, el \newif de este archivo lo vuelve a poner en
%  falso y la versión docente sale sin ninguna solución.
% =====================================================================

% --- Tipografías: la familia Source, que es la del sitio -------------
% Source Serif 4 para títulos (igual que la web) y Source Sans para el
% texto. Inter, la sans del sitio, no existe como paquete de LaTeX; Source
% Sans es la compañera de diseño de Source Serif, así que la pareja se
% mantiene coherente.
\usepackage[T1]{fontenc}
\usepackage[utf8]{inputenc}
\usepackage[default,semibold]{sourcesanspro}
\usepackage{sourceserifpro}
\usepackage[scale=0.92]{sourcecodepro}
\usepackage[spanish,es-noquoting,es-nodecimaldot]{babel}
\usepackage{amsmath,amssymb}
\usepackage[table]{xcolor}
\usepackage{booktabs}
\usepackage{enumitem}
\usepackage{fancyhdr}
\usepackage[most]{tcolorbox}
\usepackage{lastpage}
\usepackage{multicol}
\usepackage{etoolbox}
\usepackage{microtype}
\usepackage{pgffor}

\newif\ifdocente \docentefalse
\newif\ifespacio \espaciotrue

% `sourcesanspro[default]` pone la sans como tipografía del documento y
% `sourceserifpro` deja la serif en \rmfamily. Ésa es la pareja del sitio:
% texto en sans, títulos en serif.
\newcommand{\serif}{\rmfamily}

% --- Neutros y acentos ----------------------------------------------
\definecolor{tinta}{HTML}{0F172A}       % slate-900
\definecolor{tintasuave}{HTML}{334155}  % slate-700
\definecolor{apagado}{HTML}{64748B}     % slate-500
\definecolor{borde}{HTML}{E2E8F0}       % slate-200
\definecolor{papel}{HTML}{F8FAFC}       % slate-50

% --- Los bloques temáticos del Aula, con su color -------------------
\definecolor{preambulo}{HTML}{475569}     % pizarra
\definecolor{fundamentos}{HTML}{2563EB}   % azul     2.1 y 2.2
\definecolor{calculo}{HTML}{4F46E5}       % índigo   2.3 y 2.4
\definecolor{climax}{HTML}{D97706}        % ámbar    2.6
\definecolor{distrib}{HTML}{059669}       % esmeralda 2.7 a 2.9

% Color activo del apartado en curso. Cada apartado lo fija con \bloque.
\colorlet{acento}{fundamentos}
\newcommand{\bloque}[1]{\colorlet{acento}{#1}}

% --- Niveles de dificultad ------------------------------------------
\definecolor{nivelcero}{HTML}{64748B}
\definecolor{niveluno}{HTML}{2563EB}
\definecolor{niveldos}{HTML}{4F46E5}
\definecolor{niveltres}{HTML}{C2700A}
\definecolor{nivelcuatro}{HTML}{047857}

\newcommand{\nombrenivel}[1]{%
  \ifcase#1 Reconocer\or Contar\or Una fórmula\or Datos reales\or Decidir\fi}
\newcommand{\colornivel}[1]{%
  \ifcase#1 nivelcero\or niveluno\or niveldos\or niveltres\or nivelcuatro\fi}

\color{tinta}

% --- Encabezado y pie ------------------------------------------------
\pagestyle{fancy}
\fancyhf{}
\renewcommand{\headrulewidth}{0pt}
\renewcommand{\footrulewidth}{0pt}
\fancyhead[L]{\footnotesize\color{apagado}\MakeUppercase{\sffamily Psicoestadística Inferencial}}
\fancyhead[R]{\footnotesize\color{apagado}\sffamily Unidad 2 · Probabilidad}
\fancyfoot[C]{\footnotesize\color{apagado}\sffamily\thepage\,/\,\pageref{LastPage}}

\setlength{\parindent}{0pt}
\setlength{\parskip}{0.55em}

% =====================================================================
%  Encabezado de apartado — barra de color del bloque, como en la Aula
% =====================================================================
\newcommand{\apartado}[3]{%
  \clearpage
  \bloque{#3}%
  {\color{acento}\rule{\linewidth}{2.5pt}}\par
  \vspace{0.2em}
  {\footnotesize\sffamily\bfseries\color{acento}\MakeUppercase{Apartado #1}}\par
  \vspace{-0.3em}
  {\serif\fontsize{25}{28}\selectfont\bfseries #2}\par
  \vspace{0.4em}}

% --- Bajada del apartado (la frase que va bajo el título) ------------
\newcommand{\bajada}[1]{%
  {\color{tintasuave}\itshape #1}\par\vspace{0.5em}}

% =====================================================================
%  Componentes narrativos del Aula
% =====================================================================

% \definicion{término}{texto} — el recuadro con barra lateral de acento
\newcommand{\definicion}[2]{%
  \begin{tcolorbox}[enhanced, breakable, colback=white, colframe=borde,
    boxrule=0.5pt, arc=2pt, left=12pt, right=10pt, top=8pt, bottom=8pt,
    borderline west={2.5pt}{0pt}{acento}]
    {\footnotesize\sffamily\bfseries\color{acento}\MakeUppercase{Definición}}\par
    \vspace{0.15em}
    {\serif\bfseries\large #1}\par
    \vspace{0.25em}
    #2
  \end{tcolorbox}}

% \minihistoria{título}{texto} — la analogía que hace visible la idea
\newcommand{\minihistoria}[2]{%
  \begin{tcolorbox}[enhanced, breakable, colback=papel, colframe=papel,
    boxrule=0pt, arc=3pt, left=10pt, right=10pt, top=8pt, bottom=8pt]
    {\footnotesize\sffamily\bfseries\color{apagado}\MakeUppercase{#1}}\par
    \vspace{0.15em}\small #2
  \end{tcolorbox}}

% \trampa{error}{por qué}{corrección} — el error que se comete siempre
\definecolor{alertafondo}{HTML}{FEF2F2}
\definecolor{alertaborde}{HTML}{DC2626}
\newcommand{\trampa}[3]{%
  \begin{tcolorbox}[enhanced, breakable, colback=alertafondo, colframe=alertafondo,
    boxrule=0pt, arc=3pt, left=12pt, right=10pt, top=8pt, bottom=8pt,
    borderline west={2.5pt}{0pt}{alertaborde}]
    {\footnotesize\sffamily\bfseries\color{alertaborde}\MakeUppercase{La trampa}}\par
    \vspace{0.15em}
    {\small\textbf{El error:} #1\par
     \textbf{Por qué se comete:} #2\par
     \textbf{Lo correcto:} #3}
  \end{tcolorbox}}

% \hilo{texto} — el conector narrativo entre dos momentos
\newcommand{\hilo}[1]{%
  \par\vspace{0.3em}
  {\color{tintasuave}\itshape\hspace*{0.6em}\rule[0.3em]{1.6em}{0.5pt}\hspace{0.5em}#1}
  \par\vspace{0.3em}}

% \cierre{texto} — lo que quedó demostrado
\newcommand{\cierre}[1]{%
  \begin{tcolorbox}[enhanced, breakable, colback=white, colframe=borde,
    boxrule=0.5pt, arc=2pt, left=10pt, right=10pt, top=8pt, bottom=8pt]
    {\footnotesize\sffamily\bfseries\color{apagado}\MakeUppercase{Lo que acabás de ver}}\par
    \vspace{0.15em}\small #1
  \end{tcolorbox}}

% \formulaanotada{fórmula}{lista de partes} — la fórmula con sus partes
% explicadas debajo, que es como se presenta en la Aula.
\newcommand{\formulaanotada}[2]{%
  \begin{tcolorbox}[enhanced, breakable, colback=papel, colframe=borde,
    boxrule=0.5pt, arc=2pt, left=10pt, right=10pt, top=10pt, bottom=8pt]
    \centering #1\par
    \vspace{0.5em}
    \raggedright\footnotesize
    \begin{description}[leftmargin=0pt, labelsep=0.5em, itemsep=0.15em,
                        style=sameline, font=\normalfont\bfseries\color{acento}]
      #2
    \end{description}
  \end{tcolorbox}}

% =====================================================================
%  Ejercicios
% =====================================================================
\newcounter{numejercicio}
\setcounter{numejercicio}{0}

% Etiqueta del tipo de ejercicio. Los cuatro tipos nuevos existen porque la
% auditoría mostró que el 62% de los enunciados sólo pedía calcular: mirar la
% fórmula de arriba y reemplazar números. Nombrar el tipo en el margen le
% avisa al estudiante qué se le está pidiendo, que no siempre es «resolvé».
\newcommand{\tipoejercicio}[1]{%
  \ifstrempty{#1}{}{%
    \raisebox{0.1em}{\footnotesize\sffamily\color{apagado}%
      \fbox{\vphantom{Ay}\,#1\,}}}}

% \ejercicio{nivel}{tipo}{enunciado}
\newcommand{\ejercicio}[3]{%
  \stepcounter{numejercicio}%
  \par\vspace{1.2em}%
  \noindent
  {\serif\bfseries\large\color{\colornivel{#1}}\thenumejercicio.}\quad
  \raisebox{0.12em}{\footnotesize\sffamily\bfseries
    \colorbox{\colornivel{#1}}{\textcolor{white}{\,N#1 · \nombrenivel{#1}\,}}}%
  \hspace{0.35em}\tipoejercicio{#3}%
  \par\vspace{0.3em}%
  \noindent #2\par}

\newcommand{\espacioresolver}[1][3.2cm]{\ifespacio\vspace{#1}\fi}

% Renglones para escribir, en vez de un hueco vacío: invita a desarrollar y
% no sólo a poner el resultado.
\newcommand{\renglones}[1][4]{%
  \ifespacio
    \par\vspace{0.5em}
    \foreach \n in {1,...,#1}{\color{borde}\rule{\linewidth}{0.4pt}\par\vspace{0.85em}}
  \fi}

\newcommand{\solucion}[1]{%
  \ifdocente
    \par\vspace{0.35em}
    \begin{tcolorbox}[enhanced, breakable, colback=papel, colframe=borde,
      boxrule=0.5pt, arc=2pt, left=10pt, right=10pt, top=7pt, bottom=7pt]
      {\footnotesize\sffamily\bfseries\color{apagado}\MakeUppercase{Solución}}\par
      \vspace{0.1em}\small #1
    \end{tcolorbox}
  \fi}

% Nota de conducción: qué mirar mientras lo resuelven. Sólo docente.
\newcommand{\enclase}[1]{%
  \ifdocente
    \par\vspace{-0.2em}
    {\footnotesize\color{apagado}\textbf{En clase:} \itshape #1\par}
  \fi}

% La clave de respuestas se acumula en memoria: escribirla a un archivo
% auxiliar expande el contenido y rompe cualquier fórmula matemática.
\newcommand{\claves}{}
\newcommand{\respuesta}[1]{%
  \edef\numactual{\thenumejercicio}%
  \expandafter\gappto\expandafter\claves\expandafter{%
    \expandafter\item\expandafter[\numactual.]}%
  \gappto\claves{#1}}

\newcommand{\imprimirclave}{%
  \ifdocente\else
    \clearpage
    {\color{preambulo}\rule{\linewidth}{2.5pt}}\par
    \vspace{0.2em}
    {\serif\fontsize{25}{28}\selectfont\bfseries Respuestas}\par
    \vspace{0.3em}
    {\color{tintasuave}\itshape Sólo el resultado. Si tu número no coincide, el
    error está en algún paso intermedio: vale mucho más encontrarlo vos que
    leer la solución.}\par
    \vspace{0.9em}
    \begin{multicols}{2}\raggedright\small
    \begin{description}[leftmargin=2.8em, style=sameline, itemsep=0.4em]
      \claves
    \end{description}
    \end{multicols}
  \fi}
:
%     \docentetrue / \docentefalse    la solución paso a paso
%     \espaciotrue / \espaciofalse    espacio en blanco para resolver
%
%  Se declaran acá y no en cada documento: si el documento hace su propio
%  \newif antes del \input, el \newif de este archivo lo vuelve a poner en
%  falso y la versión docente sale sin ninguna solución.
% =====================================================================

% --- Tipografías: la familia Source, que es la del sitio -------------
% Source Serif 4 para títulos (igual que la web) y Source Sans para el
% texto. Inter, la sans del sitio, no existe como paquete de LaTeX; Source
% Sans es la compañera de diseño de Source Serif, así que la pareja se
% mantiene coherente.
\usepackage[T1]{fontenc}
\usepackage[utf8]{inputenc}
\usepackage[default,semibold]{sourcesanspro}
\usepackage{sourceserifpro}
\usepackage[scale=0.92]{sourcecodepro}
\usepackage[spanish,es-noquoting,es-nodecimaldot]{babel}
\usepackage{amsmath,amssymb}
\usepackage[table]{xcolor}
\usepackage{booktabs}
\usepackage{enumitem}
\usepackage{fancyhdr}
\usepackage[most]{tcolorbox}
\usepackage{lastpage}
\usepackage{multicol}
\usepackage{etoolbox}
\usepackage{microtype}
\usepackage{pgffor}

\newif\ifdocente \docentefalse
\newif\ifespacio \espaciotrue

% `sourcesanspro[default]` pone la sans como tipografía del documento y
% `sourceserifpro` deja la serif en \rmfamily. Ésa es la pareja del sitio:
% texto en sans, títulos en serif.
\newcommand{\serif}{\rmfamily}

% --- Neutros y acentos ----------------------------------------------
\definecolor{tinta}{HTML}{0F172A}       % slate-900
\definecolor{tintasuave}{HTML}{334155}  % slate-700
\definecolor{apagado}{HTML}{64748B}     % slate-500
\definecolor{borde}{HTML}{E2E8F0}       % slate-200
\definecolor{papel}{HTML}{F8FAFC}       % slate-50

% --- Los bloques temáticos del Aula, con su color -------------------
\definecolor{preambulo}{HTML}{475569}     % pizarra
\definecolor{fundamentos}{HTML}{2563EB}   % azul     2.1 y 2.2
\definecolor{calculo}{HTML}{4F46E5}       % índigo   2.3 y 2.4
\definecolor{climax}{HTML}{D97706}        % ámbar    2.6
\definecolor{distrib}{HTML}{059669}       % esmeralda 2.7 a 2.9

% Color activo del apartado en curso. Cada apartado lo fija con \bloque.
\colorlet{acento}{fundamentos}
\newcommand{\bloque}[1]{\colorlet{acento}{#1}}

% --- Niveles de dificultad ------------------------------------------
\definecolor{nivelcero}{HTML}{64748B}
\definecolor{niveluno}{HTML}{2563EB}
\definecolor{niveldos}{HTML}{4F46E5}
\definecolor{niveltres}{HTML}{C2700A}
\definecolor{nivelcuatro}{HTML}{047857}

\newcommand{\nombrenivel}[1]{%
  \ifcase#1 Reconocer\or Contar\or Una fórmula\or Datos reales\or Decidir\fi}
\newcommand{\colornivel}[1]{%
  \ifcase#1 nivelcero\or niveluno\or niveldos\or niveltres\or nivelcuatro\fi}

\color{tinta}

% --- Encabezado y pie ------------------------------------------------
\pagestyle{fancy}
\fancyhf{}
\renewcommand{\headrulewidth}{0pt}
\renewcommand{\footrulewidth}{0pt}
\fancyhead[L]{\footnotesize\color{apagado}\MakeUppercase{\sffamily Psicoestadística Inferencial}}
\fancyhead[R]{\footnotesize\color{apagado}\sffamily Unidad 2 · Probabilidad}
\fancyfoot[C]{\footnotesize\color{apagado}\sffamily\thepage\,/\,\pageref{LastPage}}

\setlength{\parindent}{0pt}
\setlength{\parskip}{0.55em}

% =====================================================================
%  Encabezado de apartado — barra de color del bloque, como en la Aula
% =====================================================================
\newcommand{\apartado}[3]{%
  \clearpage
  \bloque{#3}%
  {\color{acento}\rule{\linewidth}{2.5pt}}\par
  \vspace{0.2em}
  {\footnotesize\sffamily\bfseries\color{acento}\MakeUppercase{Apartado #1}}\par
  \vspace{-0.3em}
  {\serif\fontsize{25}{28}\selectfont\bfseries #2}\par
  \vspace{0.4em}}

% --- Bajada del apartado (la frase que va bajo el título) ------------
\newcommand{\bajada}[1]{%
  {\color{tintasuave}\itshape #1}\par\vspace{0.5em}}

% =====================================================================
%  Componentes narrativos del Aula
% =====================================================================

% \definicion{término}{texto} — el recuadro con barra lateral de acento
\newcommand{\definicion}[2]{%
  \begin{tcolorbox}[enhanced, breakable, colback=white, colframe=borde,
    boxrule=0.5pt, arc=2pt, left=12pt, right=10pt, top=8pt, bottom=8pt,
    borderline west={2.5pt}{0pt}{acento}]
    {\footnotesize\sffamily\bfseries\color{acento}\MakeUppercase{Definición}}\par
    \vspace{0.15em}
    {\serif\bfseries\large #1}\par
    \vspace{0.25em}
    #2
  \end{tcolorbox}}

% \minihistoria{título}{texto} — la analogía que hace visible la idea
\newcommand{\minihistoria}[2]{%
  \begin{tcolorbox}[enhanced, breakable, colback=papel, colframe=papel,
    boxrule=0pt, arc=3pt, left=10pt, right=10pt, top=8pt, bottom=8pt]
    {\footnotesize\sffamily\bfseries\color{apagado}\MakeUppercase{#1}}\par
    \vspace{0.15em}\small #2
  \end{tcolorbox}}

% \trampa{error}{por qué}{corrección} — el error que se comete siempre
\definecolor{alertafondo}{HTML}{FEF2F2}
\definecolor{alertaborde}{HTML}{DC2626}
\newcommand{\trampa}[3]{%
  \begin{tcolorbox}[enhanced, breakable, colback=alertafondo, colframe=alertafondo,
    boxrule=0pt, arc=3pt, left=12pt, right=10pt, top=8pt, bottom=8pt,
    borderline west={2.5pt}{0pt}{alertaborde}]
    {\footnotesize\sffamily\bfseries\color{alertaborde}\MakeUppercase{La trampa}}\par
    \vspace{0.15em}
    {\small\textbf{El error:} #1\par
     \textbf{Por qué se comete:} #2\par
     \textbf{Lo correcto:} #3}
  \end{tcolorbox}}

% \hilo{texto} — el conector narrativo entre dos momentos
\newcommand{\hilo}[1]{%
  \par\vspace{0.3em}
  {\color{tintasuave}\itshape\hspace*{0.6em}\rule[0.3em]{1.6em}{0.5pt}\hspace{0.5em}#1}
  \par\vspace{0.3em}}

% \cierre{texto} — lo que quedó demostrado
\newcommand{\cierre}[1]{%
  \begin{tcolorbox}[enhanced, breakable, colback=white, colframe=borde,
    boxrule=0.5pt, arc=2pt, left=10pt, right=10pt, top=8pt, bottom=8pt]
    {\footnotesize\sffamily\bfseries\color{apagado}\MakeUppercase{Lo que acabás de ver}}\par
    \vspace{0.15em}\small #1
  \end{tcolorbox}}

% \formulaanotada{fórmula}{lista de partes} — la fórmula con sus partes
% explicadas debajo, que es como se presenta en la Aula.
\newcommand{\formulaanotada}[2]{%
  \begin{tcolorbox}[enhanced, breakable, colback=papel, colframe=borde,
    boxrule=0.5pt, arc=2pt, left=10pt, right=10pt, top=10pt, bottom=8pt]
    \centering #1\par
    \vspace{0.5em}
    \raggedright\footnotesize
    \begin{description}[leftmargin=0pt, labelsep=0.5em, itemsep=0.15em,
                        style=sameline, font=\normalfont\bfseries\color{acento}]
      #2
    \end{description}
  \end{tcolorbox}}

% =====================================================================
%  Ejercicios
% =====================================================================
\newcounter{numejercicio}
\setcounter{numejercicio}{0}

% Etiqueta del tipo de ejercicio. Los cuatro tipos nuevos existen porque la
% auditoría mostró que el 62% de los enunciados sólo pedía calcular: mirar la
% fórmula de arriba y reemplazar números. Nombrar el tipo en el margen le
% avisa al estudiante qué se le está pidiendo, que no siempre es «resolvé».
\newcommand{\tipoejercicio}[1]{%
  \ifstrempty{#1}{}{%
    \raisebox{0.1em}{\footnotesize\sffamily\color{apagado}%
      \fbox{\vphantom{Ay}\,#1\,}}}}

% \ejercicio{nivel}{tipo}{enunciado}
\newcommand{\ejercicio}[3]{%
  \stepcounter{numejercicio}%
  \par\vspace{1.2em}%
  \noindent
  {\serif\bfseries\large\color{\colornivel{#1}}\thenumejercicio.}\quad
  \raisebox{0.12em}{\footnotesize\sffamily\bfseries
    \colorbox{\colornivel{#1}}{\textcolor{white}{\,N#1 · \nombrenivel{#1}\,}}}%
  \hspace{0.35em}\tipoejercicio{#3}%
  \par\vspace{0.3em}%
  \noindent #2\par}

\newcommand{\espacioresolver}[1][3.2cm]{\ifespacio\vspace{#1}\fi}

% Renglones para escribir, en vez de un hueco vacío: invita a desarrollar y
% no sólo a poner el resultado.
\newcommand{\renglones}[1][4]{%
  \ifespacio
    \par\vspace{0.5em}
    \foreach \n in {1,...,#1}{\color{borde}\rule{\linewidth}{0.4pt}\par\vspace{0.85em}}
  \fi}

\newcommand{\solucion}[1]{%
  \ifdocente
    \par\vspace{0.35em}
    \begin{tcolorbox}[enhanced, breakable, colback=papel, colframe=borde,
      boxrule=0.5pt, arc=2pt, left=10pt, right=10pt, top=7pt, bottom=7pt]
      {\footnotesize\sffamily\bfseries\color{apagado}\MakeUppercase{Solución}}\par
      \vspace{0.1em}\small #1
    \end{tcolorbox}
  \fi}

% Nota de conducción: qué mirar mientras lo resuelven. Sólo docente.
\newcommand{\enclase}[1]{%
  \ifdocente
    \par\vspace{-0.2em}
    {\footnotesize\color{apagado}\textbf{En clase:} \itshape #1\par}
  \fi}

% La clave de respuestas se acumula en memoria: escribirla a un archivo
% auxiliar expande el contenido y rompe cualquier fórmula matemática.
\newcommand{\claves}{}
\newcommand{\respuesta}[1]{%
  \edef\numactual{\thenumejercicio}%
  \expandafter\gappto\expandafter\claves\expandafter{%
    \expandafter\item\expandafter[\numactual.]}%
  \gappto\claves{#1}}

\newcommand{\imprimirclave}{%
  \ifdocente\else
    \clearpage
    {\color{preambulo}\rule{\linewidth}{2.5pt}}\par
    \vspace{0.2em}
    {\serif\fontsize{25}{28}\selectfont\bfseries Respuestas}\par
    \vspace{0.3em}
    {\color{tintasuave}\itshape Sólo el resultado. Si tu número no coincide, el
    error está en algún paso intermedio: vale mucho más encontrarlo vos que
    leer la solución.}\par
    \vspace{0.9em}
    \begin{multicols}{2}\raggedright\small
    \begin{description}[leftmargin=2.8em, style=sameline, itemsep=0.4em]
      \claves
    \end{description}
    \end{multicols}
  \fi}
:
%     \docentetrue / \docentefalse    la solución paso a paso
%     \espaciotrue / \espaciofalse    espacio en blanco para resolver
%
%  Se declaran acá y no en cada documento: si el documento hace su propio
%  \newif antes del \input, el \newif de este archivo lo vuelve a poner en
%  falso y la versión docente sale sin ninguna solución.
% =====================================================================

% --- Tipografías: la familia Source, que es la del sitio -------------
% Source Serif 4 para títulos (igual que la web) y Source Sans para el
% texto. Inter, la sans del sitio, no existe como paquete de LaTeX; Source
% Sans es la compañera de diseño de Source Serif, así que la pareja se
% mantiene coherente.
\usepackage[T1]{fontenc}
\usepackage[utf8]{inputenc}
\usepackage[default,semibold]{sourcesanspro}
\usepackage{sourceserifpro}
\usepackage[scale=0.92]{sourcecodepro}
\usepackage[spanish,es-noquoting,es-nodecimaldot]{babel}
\usepackage{amsmath,amssymb}
\usepackage[table]{xcolor}
\usepackage{booktabs}
\usepackage{enumitem}
\usepackage{fancyhdr}
\usepackage[most]{tcolorbox}
\usepackage{lastpage}
\usepackage{multicol}
\usepackage{etoolbox}
\usepackage{microtype}
\usepackage{pgffor}

\newif\ifdocente \docentefalse
\newif\ifespacio \espaciotrue

% `sourcesanspro[default]` pone la sans como tipografía del documento y
% `sourceserifpro` deja la serif en \rmfamily. Ésa es la pareja del sitio:
% texto en sans, títulos en serif.
\newcommand{\serif}{\rmfamily}

% --- Neutros y acentos ----------------------------------------------
\definecolor{tinta}{HTML}{0F172A}       % slate-900
\definecolor{tintasuave}{HTML}{334155}  % slate-700
\definecolor{apagado}{HTML}{64748B}     % slate-500
\definecolor{borde}{HTML}{E2E8F0}       % slate-200
\definecolor{papel}{HTML}{F8FAFC}       % slate-50

% --- Los bloques temáticos del Aula, con su color -------------------
\definecolor{preambulo}{HTML}{475569}     % pizarra
\definecolor{fundamentos}{HTML}{2563EB}   % azul     2.1 y 2.2
\definecolor{calculo}{HTML}{4F46E5}       % índigo   2.3 y 2.4
\definecolor{climax}{HTML}{D97706}        % ámbar    2.6
\definecolor{distrib}{HTML}{059669}       % esmeralda 2.7 a 2.9

% Color activo del apartado en curso. Cada apartado lo fija con \bloque.
\colorlet{acento}{fundamentos}
\newcommand{\bloque}[1]{\colorlet{acento}{#1}}

% --- Niveles de dificultad ------------------------------------------
\definecolor{nivelcero}{HTML}{64748B}
\definecolor{niveluno}{HTML}{2563EB}
\definecolor{niveldos}{HTML}{4F46E5}
\definecolor{niveltres}{HTML}{C2700A}
\definecolor{nivelcuatro}{HTML}{047857}

\newcommand{\nombrenivel}[1]{%
  \ifcase#1 Reconocer\or Contar\or Una fórmula\or Datos reales\or Decidir\fi}
\newcommand{\colornivel}[1]{%
  \ifcase#1 nivelcero\or niveluno\or niveldos\or niveltres\or nivelcuatro\fi}

\color{tinta}

% --- Encabezado y pie ------------------------------------------------
\pagestyle{fancy}
\fancyhf{}
\renewcommand{\headrulewidth}{0pt}
\renewcommand{\footrulewidth}{0pt}
\fancyhead[L]{\footnotesize\color{apagado}\MakeUppercase{\sffamily Psicoestadística Inferencial}}
\fancyhead[R]{\footnotesize\color{apagado}\sffamily Unidad 2 · Probabilidad}
\fancyfoot[C]{\footnotesize\color{apagado}\sffamily\thepage\,/\,\pageref{LastPage}}

\setlength{\parindent}{0pt}
\setlength{\parskip}{0.55em}

% =====================================================================
%  Encabezado de apartado — barra de color del bloque, como en la Aula
% =====================================================================
\newcommand{\apartado}[3]{%
  \clearpage
  \bloque{#3}%
  {\color{acento}\rule{\linewidth}{2.5pt}}\par
  \vspace{0.2em}
  {\footnotesize\sffamily\bfseries\color{acento}\MakeUppercase{Apartado #1}}\par
  \vspace{-0.3em}
  {\serif\fontsize{25}{28}\selectfont\bfseries #2}\par
  \vspace{0.4em}}

% --- Bajada del apartado (la frase que va bajo el título) ------------
\newcommand{\bajada}[1]{%
  {\color{tintasuave}\itshape #1}\par\vspace{0.5em}}

% =====================================================================
%  Componentes narrativos del Aula
% =====================================================================

% \definicion{término}{texto} — el recuadro con barra lateral de acento
\newcommand{\definicion}[2]{%
  \begin{tcolorbox}[enhanced, breakable, colback=white, colframe=borde,
    boxrule=0.5pt, arc=2pt, left=12pt, right=10pt, top=8pt, bottom=8pt,
    borderline west={2.5pt}{0pt}{acento}]
    {\footnotesize\sffamily\bfseries\color{acento}\MakeUppercase{Definición}}\par
    \vspace{0.15em}
    {\serif\bfseries\large #1}\par
    \vspace{0.25em}
    #2
  \end{tcolorbox}}

% \minihistoria{título}{texto} — la analogía que hace visible la idea
\newcommand{\minihistoria}[2]{%
  \begin{tcolorbox}[enhanced, breakable, colback=papel, colframe=papel,
    boxrule=0pt, arc=3pt, left=10pt, right=10pt, top=8pt, bottom=8pt]
    {\footnotesize\sffamily\bfseries\color{apagado}\MakeUppercase{#1}}\par
    \vspace{0.15em}\small #2
  \end{tcolorbox}}

% \trampa{error}{por qué}{corrección} — el error que se comete siempre
\definecolor{alertafondo}{HTML}{FEF2F2}
\definecolor{alertaborde}{HTML}{DC2626}
\newcommand{\trampa}[3]{%
  \begin{tcolorbox}[enhanced, breakable, colback=alertafondo, colframe=alertafondo,
    boxrule=0pt, arc=3pt, left=12pt, right=10pt, top=8pt, bottom=8pt,
    borderline west={2.5pt}{0pt}{alertaborde}]
    {\footnotesize\sffamily\bfseries\color{alertaborde}\MakeUppercase{La trampa}}\par
    \vspace{0.15em}
    {\small\textbf{El error:} #1\par
     \textbf{Por qué se comete:} #2\par
     \textbf{Lo correcto:} #3}
  \end{tcolorbox}}

% \hilo{texto} — el conector narrativo entre dos momentos
\newcommand{\hilo}[1]{%
  \par\vspace{0.3em}
  {\color{tintasuave}\itshape\hspace*{0.6em}\rule[0.3em]{1.6em}{0.5pt}\hspace{0.5em}#1}
  \par\vspace{0.3em}}

% \cierre{texto} — lo que quedó demostrado
\newcommand{\cierre}[1]{%
  \begin{tcolorbox}[enhanced, breakable, colback=white, colframe=borde,
    boxrule=0.5pt, arc=2pt, left=10pt, right=10pt, top=8pt, bottom=8pt]
    {\footnotesize\sffamily\bfseries\color{apagado}\MakeUppercase{Lo que acabás de ver}}\par
    \vspace{0.15em}\small #1
  \end{tcolorbox}}

% \formulaanotada{fórmula}{lista de partes} — la fórmula con sus partes
% explicadas debajo, que es como se presenta en la Aula.
\newcommand{\formulaanotada}[2]{%
  \begin{tcolorbox}[enhanced, breakable, colback=papel, colframe=borde,
    boxrule=0.5pt, arc=2pt, left=10pt, right=10pt, top=10pt, bottom=8pt]
    \centering #1\par
    \vspace{0.5em}
    \raggedright\footnotesize
    \begin{description}[leftmargin=0pt, labelsep=0.5em, itemsep=0.15em,
                        style=sameline, font=\normalfont\bfseries\color{acento}]
      #2
    \end{description}
  \end{tcolorbox}}

% =====================================================================
%  Ejercicios
% =====================================================================
\newcounter{numejercicio}
\setcounter{numejercicio}{0}

% Etiqueta del tipo de ejercicio. Los cuatro tipos nuevos existen porque la
% auditoría mostró que el 62% de los enunciados sólo pedía calcular: mirar la
% fórmula de arriba y reemplazar números. Nombrar el tipo en el margen le
% avisa al estudiante qué se le está pidiendo, que no siempre es «resolvé».
\newcommand{\tipoejercicio}[1]{%
  \ifstrempty{#1}{}{%
    \raisebox{0.1em}{\footnotesize\sffamily\color{apagado}%
      \fbox{\vphantom{Ay}\,#1\,}}}}

% \ejercicio{nivel}{tipo}{enunciado}
\newcommand{\ejercicio}[3]{%
  \stepcounter{numejercicio}%
  \par\vspace{1.2em}%
  \noindent
  {\serif\bfseries\large\color{\colornivel{#1}}\thenumejercicio.}\quad
  \raisebox{0.12em}{\footnotesize\sffamily\bfseries
    \colorbox{\colornivel{#1}}{\textcolor{white}{\,N#1 · \nombrenivel{#1}\,}}}%
  \hspace{0.35em}\tipoejercicio{#3}%
  \par\vspace{0.3em}%
  \noindent #2\par}

\newcommand{\espacioresolver}[1][3.2cm]{\ifespacio\vspace{#1}\fi}

% Renglones para escribir, en vez de un hueco vacío: invita a desarrollar y
% no sólo a poner el resultado.
\newcommand{\renglones}[1][4]{%
  \ifespacio
    \par\vspace{0.5em}
    \foreach \n in {1,...,#1}{\color{borde}\rule{\linewidth}{0.4pt}\par\vspace{0.85em}}
  \fi}

\newcommand{\solucion}[1]{%
  \ifdocente
    \par\vspace{0.35em}
    \begin{tcolorbox}[enhanced, breakable, colback=papel, colframe=borde,
      boxrule=0.5pt, arc=2pt, left=10pt, right=10pt, top=7pt, bottom=7pt]
      {\footnotesize\sffamily\bfseries\color{apagado}\MakeUppercase{Solución}}\par
      \vspace{0.1em}\small #1
    \end{tcolorbox}
  \fi}

% Nota de conducción: qué mirar mientras lo resuelven. Sólo docente.
\newcommand{\enclase}[1]{%
  \ifdocente
    \par\vspace{-0.2em}
    {\footnotesize\color{apagado}\textbf{En clase:} \itshape #1\par}
  \fi}

% La clave de respuestas se acumula en memoria: escribirla a un archivo
% auxiliar expande el contenido y rompe cualquier fórmula matemática.
\newcommand{\claves}{}
\newcommand{\respuesta}[1]{%
  \edef\numactual{\thenumejercicio}%
  \expandafter\gappto\expandafter\claves\expandafter{%
    \expandafter\item\expandafter[\numactual.]}%
  \gappto\claves{#1}}

\newcommand{\imprimirclave}{%
  \ifdocente\else
    \clearpage
    {\color{preambulo}\rule{\linewidth}{2.5pt}}\par
    \vspace{0.2em}
    {\serif\fontsize{25}{28}\selectfont\bfseries Respuestas}\par
    \vspace{0.3em}
    {\color{tintasuave}\itshape Sólo el resultado. Si tu número no coincide, el
    error está en algún paso intermedio: vale mucho más encontrarlo vos que
    leer la solución.}\par
    \vspace{0.9em}
    \begin{multicols}{2}\raggedright\small
    \begin{description}[leftmargin=2.8em, style=sameline, itemsep=0.4em]
      \claves
    \end{description}
    \end{multicols}
  \fi}
:
%     \docentetrue / \docentefalse    la solución paso a paso
%     \espaciotrue / \espaciofalse    espacio en blanco para resolver
%
%  Se declaran acá y no en cada documento: si el documento hace su propio
%  \newif antes del \input, el \newif de este archivo lo vuelve a poner en
%  falso y la versión docente sale sin ninguna solución.
% =====================================================================

% --- Tipografías: la familia Source, que es la del sitio -------------
% Source Serif 4 para títulos (igual que la web) y Source Sans para el
% texto. Inter, la sans del sitio, no existe como paquete de LaTeX; Source
% Sans es la compañera de diseño de Source Serif, así que la pareja se
% mantiene coherente.
\usepackage[T1]{fontenc}
\usepackage[utf8]{inputenc}
\usepackage[default,semibold]{sourcesanspro}
\usepackage{sourceserifpro}
\usepackage[scale=0.92]{sourcecodepro}
\usepackage[spanish,es-noquoting,es-nodecimaldot]{babel}
\usepackage{amsmath,amssymb}
\usepackage[table]{xcolor}
\usepackage{booktabs}
\usepackage{enumitem}
\usepackage{fancyhdr}
\usepackage[most]{tcolorbox}
\usepackage{lastpage}
\usepackage{multicol}
\usepackage{etoolbox}
\usepackage{microtype}
\usepackage{pgffor}

\newif\ifdocente \docentefalse
\newif\ifespacio \espaciotrue

% `sourcesanspro[default]` pone la sans como tipografía del documento y
% `sourceserifpro` deja la serif en \rmfamily. Ésa es la pareja del sitio:
% texto en sans, títulos en serif.
\newcommand{\serif}{\rmfamily}

% --- Neutros y acentos ----------------------------------------------
\definecolor{tinta}{HTML}{0F172A}       % slate-900
\definecolor{tintasuave}{HTML}{334155}  % slate-700
\definecolor{apagado}{HTML}{64748B}     % slate-500
\definecolor{borde}{HTML}{E2E8F0}       % slate-200
\definecolor{papel}{HTML}{F8FAFC}       % slate-50

% --- Los bloques temáticos del Aula, con su color -------------------
\definecolor{preambulo}{HTML}{475569}     % pizarra
\definecolor{fundamentos}{HTML}{2563EB}   % azul     2.1 y 2.2
\definecolor{calculo}{HTML}{4F46E5}       % índigo   2.3 y 2.4
\definecolor{climax}{HTML}{D97706}        % ámbar    2.6
\definecolor{distrib}{HTML}{059669}       % esmeralda 2.7 a 2.9

% Color activo del apartado en curso. Cada apartado lo fija con \bloque.
\colorlet{acento}{fundamentos}
\newcommand{\bloque}[1]{\colorlet{acento}{#1}}

% --- Niveles de dificultad ------------------------------------------
\definecolor{nivelcero}{HTML}{64748B}
\definecolor{niveluno}{HTML}{2563EB}
\definecolor{niveldos}{HTML}{4F46E5}
\definecolor{niveltres}{HTML}{C2700A}
\definecolor{nivelcuatro}{HTML}{047857}

\newcommand{\nombrenivel}[1]{%
  \ifcase#1 Reconocer\or Contar\or Una fórmula\or Datos reales\or Decidir\fi}
\newcommand{\colornivel}[1]{%
  \ifcase#1 nivelcero\or niveluno\or niveldos\or niveltres\or nivelcuatro\fi}

\color{tinta}

% --- Encabezado y pie ------------------------------------------------
\pagestyle{fancy}
\fancyhf{}
\renewcommand{\headrulewidth}{0pt}
\renewcommand{\footrulewidth}{0pt}
\fancyhead[L]{\footnotesize\color{apagado}\MakeUppercase{\sffamily Psicoestadística Inferencial}}
\fancyhead[R]{\footnotesize\color{apagado}\sffamily Unidad 2 · Probabilidad}
\fancyfoot[C]{\footnotesize\color{apagado}\sffamily\thepage\,/\,\pageref{LastPage}}

\setlength{\parindent}{0pt}
\setlength{\parskip}{0.55em}

% =====================================================================
%  Encabezado de apartado — barra de color del bloque, como en la Aula
% =====================================================================
\newcommand{\apartado}[3]{%
  \clearpage
  \bloque{#3}%
  {\color{acento}\rule{\linewidth}{2.5pt}}\par
  \vspace{0.2em}
  {\footnotesize\sffamily\bfseries\color{acento}\MakeUppercase{Apartado #1}}\par
  \vspace{-0.3em}
  {\serif\fontsize{25}{28}\selectfont\bfseries #2}\par
  \vspace{0.4em}}

% --- Bajada del apartado (la frase que va bajo el título) ------------
\newcommand{\bajada}[1]{%
  {\color{tintasuave}\itshape #1}\par\vspace{0.5em}}

% =====================================================================
%  Componentes narrativos del Aula
% =====================================================================

% \definicion{término}{texto} — el recuadro con barra lateral de acento
\newcommand{\definicion}[2]{%
  \begin{tcolorbox}[enhanced, breakable, colback=white, colframe=borde,
    boxrule=0.5pt, arc=2pt, left=12pt, right=10pt, top=8pt, bottom=8pt,
    borderline west={2.5pt}{0pt}{acento}]
    {\footnotesize\sffamily\bfseries\color{acento}\MakeUppercase{Definición}}\par
    \vspace{0.15em}
    {\serif\bfseries\large #1}\par
    \vspace{0.25em}
    #2
  \end{tcolorbox}}

% \minihistoria{título}{texto} — la analogía que hace visible la idea
\newcommand{\minihistoria}[2]{%
  \begin{tcolorbox}[enhanced, breakable, colback=papel, colframe=papel,
    boxrule=0pt, arc=3pt, left=10pt, right=10pt, top=8pt, bottom=8pt]
    {\footnotesize\sffamily\bfseries\color{apagado}\MakeUppercase{#1}}\par
    \vspace{0.15em}\small #2
  \end{tcolorbox}}

% \trampa{error}{por qué}{corrección} — el error que se comete siempre
\definecolor{alertafondo}{HTML}{FEF2F2}
\definecolor{alertaborde}{HTML}{DC2626}
\newcommand{\trampa}[3]{%
  \begin{tcolorbox}[enhanced, breakable, colback=alertafondo, colframe=alertafondo,
    boxrule=0pt, arc=3pt, left=12pt, right=10pt, top=8pt, bottom=8pt,
    borderline west={2.5pt}{0pt}{alertaborde}]
    {\footnotesize\sffamily\bfseries\color{alertaborde}\MakeUppercase{La trampa}}\par
    \vspace{0.15em}
    {\small\textbf{El error:} #1\par
     \textbf{Por qué se comete:} #2\par
     \textbf{Lo correcto:} #3}
  \end{tcolorbox}}

% \hilo{texto} — el conector narrativo entre dos momentos
\newcommand{\hilo}[1]{%
  \par\vspace{0.3em}
  {\color{tintasuave}\itshape\hspace*{0.6em}\rule[0.3em]{1.6em}{0.5pt}\hspace{0.5em}#1}
  \par\vspace{0.3em}}

% \cierre{texto} — lo que quedó demostrado
\newcommand{\cierre}[1]{%
  \begin{tcolorbox}[enhanced, breakable, colback=white, colframe=borde,
    boxrule=0.5pt, arc=2pt, left=10pt, right=10pt, top=8pt, bottom=8pt]
    {\footnotesize\sffamily\bfseries\color{apagado}\MakeUppercase{Lo que acabás de ver}}\par
    \vspace{0.15em}\small #1
  \end{tcolorbox}}

% \formulaanotada{fórmula}{lista de partes} — la fórmula con sus partes
% explicadas debajo, que es como se presenta en la Aula.
\newcommand{\formulaanotada}[2]{%
  \begin{tcolorbox}[enhanced, breakable, colback=papel, colframe=borde,
    boxrule=0.5pt, arc=2pt, left=10pt, right=10pt, top=10pt, bottom=8pt]
    \centering #1\par
    \vspace{0.5em}
    \raggedright\footnotesize
    \begin{description}[leftmargin=0pt, labelsep=0.5em, itemsep=0.15em,
                        style=sameline, font=\normalfont\bfseries\color{acento}]
      #2
    \end{description}
  \end{tcolorbox}}

% =====================================================================
%  Ejercicios
% =====================================================================
\newcounter{numejercicio}
\setcounter{numejercicio}{0}

% Etiqueta del tipo de ejercicio. Los cuatro tipos nuevos existen porque la
% auditoría mostró que el 62% de los enunciados sólo pedía calcular: mirar la
% fórmula de arriba y reemplazar números. Nombrar el tipo en el margen le
% avisa al estudiante qué se le está pidiendo, que no siempre es «resolvé».
\newcommand{\tipoejercicio}[1]{%
  \ifstrempty{#1}{}{%
    \raisebox{0.1em}{\footnotesize\sffamily\color{apagado}%
      \fbox{\vphantom{Ay}\,#1\,}}}}

% \ejercicio{nivel}{tipo}{enunciado}
\newcommand{\ejercicio}[3]{%
  \stepcounter{numejercicio}%
  \par\vspace{1.2em}%
  \noindent
  {\serif\bfseries\large\color{\colornivel{#1}}\thenumejercicio.}\quad
  \raisebox{0.12em}{\footnotesize\sffamily\bfseries
    \colorbox{\colornivel{#1}}{\textcolor{white}{\,N#1 · \nombrenivel{#1}\,}}}%
  \hspace{0.35em}\tipoejercicio{#3}%
  \par\vspace{0.3em}%
  \noindent #2\par}

\newcommand{\espacioresolver}[1][3.2cm]{\ifespacio\vspace{#1}\fi}

% Renglones para escribir, en vez de un hueco vacío: invita a desarrollar y
% no sólo a poner el resultado.
\newcommand{\renglones}[1][4]{%
  \ifespacio
    \par\vspace{0.5em}
    \foreach \n in {1,...,#1}{\color{borde}\rule{\linewidth}{0.4pt}\par\vspace{0.85em}}
  \fi}

\newcommand{\solucion}[1]{%
  \ifdocente
    \par\vspace{0.35em}
    \begin{tcolorbox}[enhanced, breakable, colback=papel, colframe=borde,
      boxrule=0.5pt, arc=2pt, left=10pt, right=10pt, top=7pt, bottom=7pt]
      {\footnotesize\sffamily\bfseries\color{apagado}\MakeUppercase{Solución}}\par
      \vspace{0.1em}\small #1
    \end{tcolorbox}
  \fi}

% Nota de conducción: qué mirar mientras lo resuelven. Sólo docente.
\newcommand{\enclase}[1]{%
  \ifdocente
    \par\vspace{-0.2em}
    {\footnotesize\color{apagado}\textbf{En clase:} \itshape #1\par}
  \fi}

% La clave de respuestas se acumula en memoria: escribirla a un archivo
% auxiliar expande el contenido y rompe cualquier fórmula matemática.
\newcommand{\claves}{}
\newcommand{\respuesta}[1]{%
  \edef\numactual{\thenumejercicio}%
  \expandafter\gappto\expandafter\claves\expandafter{%
    \expandafter\item\expandafter[\numactual.]}%
  \gappto\claves{#1}}

\newcommand{\imprimirclave}{%
  \ifdocente\else
    \clearpage
    {\color{preambulo}\rule{\linewidth}{2.5pt}}\par
    \vspace{0.2em}
    {\serif\fontsize{25}{28}\selectfont\bfseries Respuestas}\par
    \vspace{0.3em}
    {\color{tintasuave}\itshape Sólo el resultado. Si tu número no coincide, el
    error está en algún paso intermedio: vale mucho más encontrarlo vos que
    leer la solución.}\par
    \vspace{0.9em}
    \begin{multicols}{2}\raggedright\small
    \begin{description}[leftmargin=2.8em, style=sameline, itemsep=0.4em]
      \claves
    \end{description}
    \end{multicols}
  \fi}
